\section{Types and random matchings}\label{sec:random}

We construct the outgoing label sets first, so that every choice of edges
has the parity properties needed for the algebraic argument.  We then
choose the edges by random matchings.  The estimates in this section show
that we may condition on logarithmic girth while retaining useful bounds
for prescribed edges, and that the resulting components have logarithmic
diameter.  The construction and the use of conditioned matchings adapt
the method of \cite[Section ``Prescribed types and random matchings'']{OpenAI2026ZD};
the label counts and estimates below are for the types used here.

\subsection{Prescribing the outgoing labels}

Fix
\[
 q=128,\qquad v=q^2+q+1=16513,\qquad p=\frac{q+1}{v}.
\]
Let \(\mathcal P\) and \(\mathcal L\) be the points and lines of the
projective plane over \(\mathbb F_q\).  Concretely, these are the
one-dimensional and two-dimensional subspaces of \(\mathbb F_q^3\),
with incidence given by containment.  Counting nonzero vectors gives
\(v\) points, and duality gives \(v\) lines.  Every line contains
\(q+1\) points, and every point lies on \(q+1\) lines.  Two distinct
points span a unique line, and two distinct lines intersect in exactly
one point.

Adjoin a set \(\mathcal E\) of seven extra letters to \(\mathcal P\),
and put \(T=\mathcal P\sqcup\mathcal E\).  Pair each extra letter
with a distinct ordinary letter.  Since \(v-7\) is even, the remaining
ordinary letters can be paired among themselves.  These pairs define
the fixed-point-free involution \(t\mapsto\bar t\) on \(T\).
As before, \(F\) is the rose with one edge for each inverse pair.

Identify \(\mathcal E\) with the seven points of the projective plane
over \(\mathbb F_2\), and let \(\mathcal D\) consist of the complements
of its seven lines.  Each \(D\in\mathcal D\) has four elements, and
distinct members of \(\mathcal D\) intersect in two elements.  Each
extra letter belongs to four members of \(\mathcal D\): three Fano
lines contain the letter and four do not.  Each pair of distinct extra
letters belongs to two members of \(\mathcal D\), since the number of
Fano lines containing at least one of them is \(3+3-1=5\).

Let \(m\) tend to infinity through positive integers congruent to
\(1\) modulo \(4\), and set \(n=vm\).  All asymptotic statements
below use this sequence.  Partition sets \(A\) and \(B\) into classes
indexed by \(\mathcal L\), with \(m\) vertices in each class of
\(A\) and \(m-1\) vertices in each class of \(B\).  Thus
\[
 |A|=n,\qquad |B|=n-v.
\]
The ordinary part of the outgoing label set \(S_x\) at a vertex
\(x\) is its indexing line.  Choose a distinguished vertex
\(x_A\in A\) and give it all seven extra letters.  For every
\(D\in\mathcal D\), prescribe extra part exactly \(D\) at
\begin{equation}\label{eq:extra-counts}
 a_m=\frac{(q+1)m-1}{4}\quad\hbox{vertices of }A,
 \qquad
 b_m=\frac{(q+1)(m-1)}4\quad\hbox{vertices of }B.
\end{equation}
These are integers because \(q+1\equiv m\equiv1\pmod4\).
Every remaining vertex has empty extra part.

For each \(D\), distribute its prescribed vertices among the \(v\)
line classes with counts differing by at most one.  These distributions
can be realized simultaneously on disjoint vertices.  Indeed, in each
line class the total demand in \(A\), including \(x_A\) if applicable,
is at most \(7\lceil a_m/v\rceil+1\), and in \(B\) it is at most
\(7\lceil b_m/v\rceil\).  Both bounds are
\[
 \frac{7(q+1)}{4v}m+O(1),
\]
whose coefficient of \(m\) is strictly less than one.  For all
sufficiently large \(m\), each bound is smaller than \(m-1\).
Fix one such assignment of the sets \(S_x\), and also fix any root
\(x_B\in B\), before choosing edges.

\begin{proposition}[Balanced types and parity]\label{prop:types}
For all sufficiently large \(m\equiv1\pmod4\), the prescribed sets
\(S_x\subseteq T\), \(x\in A\sqcup B\), satisfy the following
identities.  For every \(Y\in\{A,B\}\) and \(t\in T\),
\begin{equation}\label{eq:label-balance}
 \bigl|\{x\in Y:t\in S_x\}\bigr|=|Y|p.
\end{equation}
Their intersection cardinalities satisfy
\begin{equation}\label{eq:type-parity}
 \begin{aligned}
 |S_x\cap S_{x'}|&\equiv1\pmod2
 &&\bigl(x,x'\in A,\ (x,x')\ne(x_A,x_A)\bigr),\\
 |S_{x_A}|&\equiv0\pmod2,\\
 |S_x\cap S_y|&\equiv1\pmod2
 &&\bigl(x\in A,\ y\in B\bigr).
 \end{aligned}
\end{equation}
Moreover, \(129\le |S_x|\le136\) at every vertex.
\end{proposition}

\begin{proof}
An ordinary letter lies on \(q+1\) lines, so occurs at exactly
\((q+1)m=np\) vertices of \(A\) and \((q+1)(m-1)=(n-v)p\)
vertices of \(B\).  An extra letter belongs to four members of
\(\mathcal D\), and therefore occurs \(4a_m+1=(q+1)m\) times in
\(A\) and \(4b_m=(q+1)(m-1)\) times in \(B\).
This proves \eqref{eq:label-balance}.

The intersection of two ordinary parts has size \(q+1=129\) when
the lines agree and size one otherwise.  For the extra parts, all
intersections have even size except the intersection of the full
seven-element set with itself.  In fact, the possible even sizes are
zero, two, and four.  Only \(x_A\) has the full extra part, giving
exactly the exception in \eqref{eq:type-parity}.  Finally, the possible
sizes of \(S_x\) are \(129\), \(133\), and \(136\).
\end{proof}

For \(Y\in\{A,B\}\), write
\(Y_t=\{x\in Y:t\in S_x\}\).  For each inverse pair, choose a
bijection \(Y_t\to Y_{\bar t}\).  A matched pair \((x,y)\) gives
one undirected edge, oriented from \(x\) to \(y\) with label \(t\)
and in the reverse direction with label \(\bar t\).  Overlap between
\(Y_t\) and \(Y_{\bar t}\) is allowed, so loops may occur at this
stage.  More precisely, a match joins the label slots \((x,t)\) and
\((y,\bar t)\); these remain distinct when \(x=y\).
Denote the resulting graphs by \(\Gamma_A\) and \(\Gamma_B\),
and put \(\Gamma=\Gamma_A\sqcup\Gamma_B\).  Every outgoing label
at a vertex occurs exactly once, so the maps to \(F\) are immersions.
Degrees count half-edges, including both half-edges of a loop, and
are bounded by
\begin{equation}\label{eq:degree-bound}
 129\le\deg(x)\le d_*:=136.
\end{equation}

\subsection{Turn frequencies and decay of word weights}

At an interior vertex of a path with consecutive labels \(t,u\), the
outgoing labels used by the path are \(\bar t,u\).  For a reduced
turn, meaning \(u\ne\bar t\), define
\begin{equation}\label{eq:turn-weights}
 w(t,u)=
 \begin{cases}
 (q+1)^{-1},&\bar t,u\in\mathcal P,\\
 p,&\text{one of }\bar t,u\text{ is in }\mathcal P
          \text{ and the other is in }\mathcal E,\\
 1/2,&\bar t,u\in\mathcal E.
 \end{cases}
 \qquad w(t,\bar t)=0.
\end{equation}
For either \(Y=A\) or \(Y=B\), uniformly over reduced turns,
\begin{equation}\label{eq:turn-count}
 \bigl|\{x\in Y:\bar t,u\in S_x\}\bigr|
       =np\,w(t,u)+O(1).
\end{equation}
Here and below constants in \(O(1)\) do not depend on \(m\).
To verify the formula, two distinct ordinary letters occur together
on just one line, giving \(m\) or \(m-1\) choices.  An extra letter
occurs in each line class \(4a_m/v+O(1)=pm+O(1)\) times in \(A\),
including the possible distinguished vertex in the error, and
\(4b_m/v+O(1)=p(m-1)+O(1)\) times in \(B\).  Summing over the
\(q+1\) lines containing a specified ordinary letter gives
\(np^2+O(1)\).  Finally, two distinct extra letters occur together
\(2a_m+1=np/2+1/2\) times in \(A\) and
\(2b_m=(n-v)p/2\) times in \(B\).  These are precisely the three
cases of \eqref{eq:turn-count}.

For a word \(W=t_1\cdots t_h\) of positive length, put
\begin{equation}\label{eq:word-weight}
 P(W)=\prod_{i=1}^{h-1}w(t_i,t_{i+1}),
\end{equation}
with the empty product equal to one.  This weight is positive exactly
when the word is reduced.  Since the rule in \eqref{eq:turn-weights}
is symmetric in \(\bar t,u\), it also gives
\[
 w(t,u)=w(\bar u,\bar t),\qquad
 P(W^{-1})=P(W),\qquad W^{-1}=\bar t_h\cdots\bar t_1.
\]
Thus reversing the orientation of a path does not change its weight.
The square of this weight has summable total mass in long word lengths,
as the next estimate shows.

\begin{lemma}[Decay of word weights]\label{lem:word-decay}
There is a constant \(\delta>0\), depending only on the prescribed
alphabet and weights, such that for all sufficiently large integers
\(h\),
\begin{equation}\label{eq:word-decay}
 \sum_{W\in T^h}P(W)^2\le \exp(-2\delta h).
\end{equation}
\end{lemma}

\begin{proof}
Let \(M\) be the nonnegative matrix with entries
\(M_{t,u}=w(t,u)^2\).  Define a positive vector \(f\) on \(T\) by
\[
 f(t)=\begin{cases}4,&\bar t\in\mathcal E,\\1,&\bar t\in\mathcal P.
 \end{cases}
\]
Exactly seven ordinary letters have value four under \(f\); all
extra letters have value one.  If \(\bar t\) is ordinary, there are
\(v-1=q(q+1)\) allowed ordinary successors and seven extra
successors.  Allowing all seven of the exceptional ordinary letters
among those successors gives
\[
 (Mf)(t)\le
 \frac{q}{q+1}+7p^2+\frac{21}{(q+1)^2}<1=f(t).
\]
If \(\bar t\) is extra, the six allowed extra successors contribute
\(6/4\), and the ordinary successors contribute \((v+21)p^2\).
Thus
\[
 (Mf)(t)\le\frac64+(v+21)p^2<4=f(t).
\]
Both strict inequalities hold for \(q=128\).  Consequently a fixed
\(\lambda<1\) satisfies \(Mf\le\lambda f\) coordinatewise.
Writing \(\mathbf1\) for the all-ones vector and using
\(\mathbf1\le f\le4\mathbf1\), we obtain
\[
 \sum_{W\in T^h}P(W)^2
 =\mathbf1^{\mathsf T}M^{h-1}\mathbf1
 \le4|T|\lambda^{h-1}.
\]
For example, \(\delta=-\tfrac14\log\lambda>0\) makes the last
quantity at most \(\exp(-2\delta h)\) for all sufficiently large
\(h\).
\end{proof}

\subsection{Conditioning on large girth}

Initially choose all the bijections independently and uniformly, and
write \(\mathbb P\) for this probability law.  Fix \(c_0>0\) so small
that
\begin{equation}\label{eq:girth-scale}
 c_0\log(2|T|/p)<1,\qquad 2c_0\log d_*<1,
 \qquad L=\lfloor c_0\log n\rfloor.
\end{equation}
Girth is the length of a shortest simple cycle; loops and pairs of
parallel edges count as cycles of lengths one and two.  Let
\(\mathcal G_L\) be the event that \(\Gamma\) has girth at least
\(L\).  We first prove that this event is nonempty for all sufficiently
large \(n\), so that the conditioned law
\[
 \mathbb P_L(\,\cdot\,)=\mathbb P(\,\cdot\mid\mathcal G_L)
\]
is defined.

A \emph{match prescription} specifies one matched pair in one of the
bijections, on a specified side \(A\) or \(B\).  The reverse
orientation describes the same prescription.  A collection is feasible
if some tuple of bijections satisfies it.  Every bijection on side
\(Y\) has \(N_Y=|Y|p\) pairs.  Before conditioning, a feasible
collection of distinct prescriptions has probability
\(\prod (N_Y)_{e}^{-1}\), where the product runs over the bijections,
\(e\) is their respective number of prescriptions, and
\((N)_e=N(N-1)\cdots(N-e+1)\).

There are at most \(2n^k|T|^k\) choices of vertices and labels that
can describe a cycle of length \(k\).  A feasible simple cycle uses
\(k\) distinct prescriptions, even for \(k=1\) or \(k=2\), and
therefore has probability at most \((np-vp-k)^{-k}\).  If an edge is
called bad when it lies on a cycle of length less than \(L\), the
expected number of bad edges is consequently at most
\begin{equation}\label{eq:bad-edges}
 2\sum_{1\le k<L} k\left(\frac{2|T|}{p}\right)^k=o(n).
\end{equation}
Here \(np-vp-k\ge np/2\) for large \(n\), uniformly for \(k<L\),
and the last assertion follows from \eqref{eq:girth-scale}.

We use the following elementary switch.  Suppose that \(xy\) and
\(uz\) are edges on the same side, both oriented with the same label
\(t\).  Assume that \(uz\) is not bad and that the distance between
\(\{x,y\}\) and \(\{u,z\}\) exceeds \(2L\).  Distances between
different components are understood to be infinite.  Transpose the
two targets in the bijection, replacing \(xy,uz\) by \(xz,uy\).
This preserves every outgoing label set and creates no cycle of length
less than \(L\).

For the last assertion, any new short cycle must contain a new edge.
If it contains just one, deleting that edge gives an old path of
length less than \(L\) between the two distant endpoint sets, a
contradiction.  If it contains both, deleting the new edges gives
two old paths.  Either one of these crosses between the endpoint
sets, again contradicting their distance, or the paths join \(x\)
to \(y\) and \(u\) to \(z\).  In the latter case the old
\(u\)-to-\(z\) path together with \(uz\) is a cycle of length
less than \(L\), contrary to the choice of \(uz\).  This argument
also applies when \(x=y\); the corresponding residual path can have
length zero.

Choose an initial tuple with \(o(n)\) bad edges, as supplied by
\eqref{eq:bad-edges}.  Given a bad edge \(xy\), the number of edges
having an endpoint within distance \(2L\) of \(x\) or \(y\) is
\(O(d_*^{2L+2})=o(n)\), by \eqref{eq:girth-scale}.  The appropriate
bijection contains at least \(np-vp\) edges, so it has an edge
\(uz\) which is neither bad nor excluded by this distance condition.
The switch removes the bad edge \(xy\) and creates no new bad edges.
Repeating decreases the number of bad edges strictly at each step
and eventually gives a tuple in \(\mathcal G_L\).  Since every
tuple has positive probability, \(\mathbb P(\mathcal G_L)>0\).

The purpose of switching is also quantitative.  It permits an upper
bound for an additional prescribed edge without any lower bound for
\(\mathbb P(\mathcal G_L)\).

\begin{lemma}[Prescriptions under girth conditioning]\label{lem:prescriptions}
There is a nonnegative sequence
\(r_n=vp+O(d_*^{2L+2})=o(n)\) with the following property.  If
\(\mathcal S\) is a collection of \(s\) distinct prescriptions
having positive probability under \(\mathbb P_L\), and \(e\) is
a further prescription not in \(\mathcal S\), then
\begin{equation}\label{eq:one-prescription}
 \mathbb P_L(e\mid\mathcal S)
 \le\frac1{np-s-r_n}
 \qquad\text{whenever }np-s-r_n>0.
\end{equation}
In particular, for every fixed \(C_1>0\), a specified collection of
\(E\le C_1L\) distinct prescriptions satisfies the uniform bound
\begin{equation}\label{eq:short-prescriptions}
 \mathbb P_L(\text{all }E\text{ prescriptions hold})
 \le (np)^{-E}\exp(o(L)).
\end{equation}
\end{lemma}

\begin{proof}
The conditional law is uniform on the finite set of tuples satisfying
\(\mathcal G_L\) and \(\mathcal S\).  Partition this set into
\(\Omega_+\), the tuples containing \(e=xy\), and \(\Omega_-\),
those not containing it.  If \(\Omega_+\) is empty there is nothing
to prove.  Otherwise orient \(e\) with its prescribed label.
For any tuple in \(\Omega_+\), at least
\[
 np-vp-s-O(d_*^{2L+2})\ge np-s-r_n
\]
edges of its bijection are sufficiently distant from \(\{x,y\}\)
and are not in \(\mathcal S\).  Choose the constant defining \(r_n\)
uniformly large enough for this inequality.  Switching \(xy\) with
any such edge gives a tuple in \(\Omega_-\): all old edges avoid
short cycles, so the preceding switching argument preserves girth,
and neither deleted edge is prescribed by \(\mathcal S\).

The output uniquely determines the input and the chosen switch.
Indeed, if \(\phi'\colon Y_t\to Y_{\bar t}\) is the output
bijection, then \(z=\phi'(x)\) and \(u=(\phi')^{-1}(y)\).
Replacing these two matches by \(\phi(x)=y\) and \(\phi(u)=z\)
recovers the input, even when the domain and codomain overlap as
vertex sets.  Thus
\[
 |\Omega_+|(np-s-r_n)\le |\Omega_-|.
\]
Dividing by the total number of tuples proves
\eqref{eq:one-prescription}.  Incompatible additional prescriptions
have probability zero and satisfy the same bound.

Expose a specified collection of \(E\le C_1L\) edges in any order.
If a prefix is impossible its joint probability is zero; otherwise
\eqref{eq:one-prescription} applies at every step.  Hence
\[
 \mathbb P_L(\text{all }E\text{ prescriptions hold})
 \le\prod_{s=0}^{E-1}(np-s-r_n)^{-1}
 \le (np)^{-E}
       \exp\!\left(O\!\left(\frac{E(E+r_n)}n\right)\right).
\]
Since \(E=O(L)\), \(L=O(\log n)\), and \(r_n=o(n)\), the error
in the exponent is \(o(L)\), uniformly for the stated range of \(E\).
\end{proof}

\subsection{Expansion and component diameter}

We now use \eqref{eq:one-prescription} for collections with a small
linear number of edges.  The high minimum degree forces small vertex
sets to expand, even under girth conditioning.

\begin{proposition}[Component diameter]\label{prop:diameter}
There is a fixed constant \(d>0\) such that, with
\(\mathbb P_L\)-probability tending to one as \(m\to\infty\), every
connected component of \(\Gamma_A\) and \(\Gamma_B\) has diameter
at most \(dL\).
\end{proposition}

\begin{proof}
For a vertex set \(S\) in either graph, let \(N[S]\) be \(S\)
together with all its neighbors.  We first show that there is a
constant \(\rho>0\) such that, with probability tending to one,
\begin{equation}\label{eq:small-set-expansion}
 |N[S]|\ge2|S|
 \quad\text{for every }S\text{ with }1\le|S|\le\rho n
 \quad\text{on either side}.
\end{equation}
All constants used to choose \(\rho\) below are independent of \(n\).

Suppose \(|S|=k\le\rho n\) and \(|N[S]|<2k\).  Take \(\rho<1/4\),
so that \(2k\le n-v\) for sufficiently large \(n\).  We can then
enlarge \(N[S]\) to a set \(Z\) of exactly \(2k\) vertices on
the same side.  Every edge incident to \(S\) has both endpoints in
\(Z\).  The sum of the degrees in \(S\) is at least \(129k\),
and an edge contributes at most two to that sum.  Therefore \(Z\)
contains at least
\[
 r=\lceil129k/2\rceil
\]
distinct edges.  Notice that \(64.5k\le r\le65k\).

There are at most \((en/(2k))^{2k}\) choices for \(Z\).  For a
fixed \(Z\), the number of possible labeled edge prescriptions
with endpoints in \(Z\) is at most \(4|T|k^2\).  Thus the number
of sets of \(r\) such prescriptions is at most
\[
 \binom{4|T|k^2}{r}\le (C_2k)^r
\]
for a fixed constant \(C_2\); if fewer than \(r\) prescriptions
are available, the event is empty.  Reduce \(\rho\) so that
\(65\rho<p/4\).  Because \(r_n=o(n)\), all denominators in
\eqref{eq:one-prescription}, when exposing these \(r\) edges, are
at least \(np/2\) for sufficiently large \(n\).  Each collection
therefore has probability at most \((2/(np))^r\).

A union bound, first over \(Z\) and then its prescribed edges, now
bounds the probability of failure for this \(k\) on either fixed
side by
\begin{equation}\label{eq:expansion-probability}
 \left(\frac{en}{2k}\right)^{2k}
 (C_2k)^r\left(\frac2{np}\right)^r
 \le\left(C_3(k/n)^{62.5}\right)^k
\end{equation}
for a fixed \(C_3\).  The exponent \(62.5\) is
\(129/2-2\); rounding \(r\) changes only the constant because
\(64.5\le r/k\le65\).  Choose \(\rho\) smaller once more so that
\(C_3\rho^{62.5}<1/2\).  For \(k\le\sqrt n\), the right side
of \eqref{eq:expansion-probability} is at most
\((C_3n^{-31.25})^k\), whose sum tends to zero.  For
\(\sqrt n<k\le\rho n\), it is at most \(2^{-k}\), again with
sum tending to zero.  Including both sides proves
\eqref{eq:small-set-expansion} with probability \(1-o(1)\).

Assume \eqref{eq:small-set-expansion}.  Starting with any vertex,
successive closed balls double in cardinality as long as they have
at most \(\rho n\) vertices.  In particular every ball of radius
\[
 R=\lceil\log_2 n\rceil+1
\]
has more than \(\rho n\) vertices.  On a geodesic of length \(D\)
in a component, select vertices at spacings \(2R+1\).  Their
radius-\(R\) balls are pairwise disjoint, because distances between
vertices on a geodesic equal their separations along it.  All these
balls lie on one side, which has at most \(n\) vertices.  Hence
\[
 \left(\left\lfloor\frac D{2R+1}\right\rfloor+1\right)\rho n<n,
 \qquad D<\frac{2R+1}{\rho}.
\]
Since \(R=O(\log n)\) and \(L\sim c_0\log n\), a fixed
\(d\) bounds the last expression by \(dL\) for all sufficiently
large \(n\).  This proves the proposition simultaneously for all
components.
\end{proof}
